\documentclass[a4paper,oneside,10pt]{article}
\usepackage{wallpaper}
\usepackage{float}
\usepackage[utf8]{inputenc}
\usepackage{circuitikz}

\usepackage{enumitem}
\usepackage{amsmath}

\setlength{\textwidth}{16cm} %
\setlength{\textheight}{23cm} %
\setlength{\oddsidemargin}{0cm} %
\setlength{\evensidemargin}{0cm} %
\setlength{\topmargin}{0cm}

\CenterWallPaper{1.01}{bg.pdf}


\begin{document}

\noindent Relatório 1: Associação de resistores. \\
Laboratório de Circuitos Elétricos - Integral\\
Prof. Gabriel de Freitas Fernandes\\
2026.2\\
\\
\textbf{EQUIPE: Leonardo Tomasela Leal}
\\

\section{Introdução}

O objetivo deste relatório é responder as questões estabelecidas em cada experimento realizado na atividade prática 1, que envolvia a criação de três circuitos elétricos diferentes, as perguntas estabelecidas eram:
\begin{enumerate}
	\item Experimento 1
	\begin{enumerate}
		\item Qual o valor teórico da associação de resistores?
		\item Monte o circuito e meça o valor da associação de resistores.
		\item Meça a corrente da malha principal e a que passa por um dos resistores de 680 $\Omega$
		\item Meça a corrente e tensão no resistor de 100 $\Omega$.
	\end{enumerate}
	\item Experimento 2
	\begin{enumerate}
		\item Encontre a equação que descreve a associação de resistores presente na imagem.
		\item Monte o circuito e meça a resistência da associação.
	\end{enumerate}
	\item Experimento 3
	\begin{enumerate}
		\item Calcule o valor ideal para o resistor R e monte o circuito. Qual valor real foi escolhido? Maior ou menor do que o ideal?
		\item Construa o circuito com um potenciômetro. Qual foi o efeito da variação da resistência?
	\end{enumerate}
\end{enumerate}

As imagens mencionadas são reproduzidas a seguir nos respectivos subtópicos.

%\section{Pré-Laboratório (Caso se aplique)}

%Exponha todas as etapas de pré-laboratório. Caso hajam múltiplas atividades, separe cada uma da seguinte forma:

%\subsection{Atividade Pré-Lab 1(Use o nome correto)}

%\subsection{Atividade Pré-Lab 2(Use o nome correto)}circuitikz

\section{Experimento 1}

\subsection{Montagem Experimental}

Foram utilizados 2 resistores de 680 $\Omega$, com tolerância de $\pm$5\% associados em paralelo entre si, e associados em série a outro resistor de 100 $\Omega$ com tolerância de $\pm$5\%

Todos foram posicionados na protoboard nessa configuração, o terminal positivo da fonte de tensão conectado à extremidade livre da associação em paralelo dos resistores de 680 $\Omega$ e o terminal negativo foi conectado na extremidade livre do resistor de 100 $\Omega$. O terminal positivo foi configurado com uma tensão contínua de 5 V e o negativo com uma tensão de 0 V, a fonte também foi configurada com uma corrente máxima de 10 mA.
\begin{center}
\begin{circuitikz}[scale=0.9, transform shape]
  \draw
% Fonte de 5V (positivo em cima, negativo no terra)
(2.75,2) node[vcc](vcc){5V} (2.75,2)

% Barramento positivo (topo)
(2,2) -- (3.5,2)

% Resistor 1 do paralelo (ramo esquerdo)
(2,2) to[R, l=680$\Omega$] (2,0)

% Resistor 2 do paralelo (ramo direito)
(3.5,2) to[R, l=680$\Omega$] (3.5,0)

% Barramento negativo do paralelo (fundo)
(2,0) -- (3.5,0) -- (2.75,0)

% Resistor de 100Ω em série descendo para o terra
-- (2.75,-0.3) to[R, l=100$\Omega$] (2.75,-1.5)

% Terminal final conectado ao terra
-- (2.75,-1.5) node[ground]{};
\end{circuitikz}
\end{center}
%\subsection{Código}


\subsection{Resultados e Discussão}


\begin{enumerate}[label=\alph*]
	\item \textbf{Qual o valor teórico da associação de resistores?}
	
	Como podemos analisar, temos 2 resistores de 680 $\Omega$ em paralelo entre si, calculando a resistência equivalente desta associação temos:
	\begin{equation}
		\frac{1}{R_{1}}=\frac{1}{680\ \Omega} + \frac{1}{680\ \Omega} = \frac{2}{680\ \Omega}\Rightarrow R_1 = \frac{680\ \Omega}{2} \Rightarrow R_1 = 340\ \Omega
		\label{eq:r_1}
	\end{equation}
	\begin{equation}
		R_2 = 100\Omega
		\label{eq:r_2}
	\end{equation}
	Logo, temos que:
	$$R_{eq} = R_1 + R_2 \Rightarrow R_{eq} = 340\ \Omega + 100\ \Omega =  \boxed{440\ \Omega}$$
	
	
	\item \textbf{Monte o circuito e meça o valor da associação de resistores.}
	
	O circuito foi montado em laborado, a medição foi feita utilizando o multímetro de bancada em modo de medição de ohmímetro. A fonte de tensão contínua não foi conectada ao circuito. O valor medido foi de \boxed{441\ \Omega}
	
	\item \textbf{Meça a corrente da malha principal e a que passa por um dos resistores de 680 $\Omega$.}
	
	\begin{enumerate}
		
		
		\item \textbf{Corrente que passa pela malha principal:}
		
		O circuito foi aberto e reestruturado para garantir que o múltimetro fosse colocado em série com a malha principal para medir corrente contínua, a seguir fica a representação do circuito com o multímetro em modo amperímetro:
		
		\begin{center}
			\begin{circuitikz}[scale=0.9, transform shape]
				\draw
				% Fonte de 5V (positivo em cima, negativo no terra)
				(2.75,3.5) node[vcc](vcc){+5V} (2.75,3.5)
				
				% Barramento positivo (topo)
				(2,3.5) -- (3.5,3.5)
				
				% Resistor 1 do paralelo (ramo esquerdo)
				(2,3.5) to[R, l=680$\Omega$] (2,1.75)
				
				
				% Resistor 2 do paralelo (ramo direito)
				
				
				(3.5,3.5) to[R, l=680$\Omega$] (3.5,1.75)
				
				% Barramento negativo do paralelo (fundo)
				(2,1.75) -- (3.5,1.75) -- (2.75,1.75)
				
				(2.75,1.75) to[R, l=100$\Omega$] (2.75,0)
				% Resistor de 100Ω em série descendo para o terra
				-- (2.75,-0.3) to[ammeter] (2.75,-1.5)
				
				% Terminal final conectado ao terra
				-- (2.75,-1.5) node[ground]{};
			\end{circuitikz}
		\end{center}
		O valor medido pelo amperímetro foi de \boxed{12\ mA}
		
		\item \textbf{Corrente que passa pelo resistor de 680 $\Omega$:}
		
		O circuito foi aberto e reestruturado para garantir que o múltimetro fosse colocado em série com o resistor para medir corrente contínua, a seguir fica a representação do circuito com o multímetro em modo amperímetro:
		
		\begin{center}
			\begin{circuitikz}[scale=0.9, transform shape]
				\draw
				% Fonte de 5V (positivo em cima, negativo no terra)
				(2.75,3.5) node[vcc](vcc){5V} (2.75,3.5)
				
				% Barramento positivo (topo)
				(2,3.5) -- (3.5,3.5)
				
				% Resistor 1 do paralelo (ramo esquerdo)
				(2,3.5) to[R, l=680$\Omega$] (2,0)
				
				
				% Resistor 2 do paralelo (ramo direito)
				
				(3.5,3.5) to[R, l=680$\Omega$] (3.5,1.75)
				(3.5,1.75) to[ammeter] (3.5,0)
				
				% Barramento negativo do paralelo (fundo)
				(2,0) -- (3.5,0) -- (2.75,0)
				
				% Resistor de 100Ω em série descendo para o terra
				-- (2.75,-0.3) to[R, l=100$\Omega$] (2.75,-1.5)
				
				% Terminal final conectado ao terra
				-- (2.75,-1.5) node[ground]{};
			\end{circuitikz}
		\end{center}
		
		O valor medido pelo amperímetro foi de \boxed{5,51\ mA}
		
	\end{enumerate}
	
	\item \textbf{Meça a corrente e tensão no resistor de 100 $\Omega$.}
	
	\begin{enumerate}
		
		\item \textbf{Corrente no resistor de 100 $\Omega$:}
		
		Para medir a corrente no resistor de 100 $\Omega$, o circuito foi aberto e foi colocado um multímetro em modo amperímetro entre o resistor de 100 $\Omega$ e os dois resistores de 680 $\Omega$, de forma que o circuito resultante fique da seguinte forma:
		
		
		\begin{center}
			\begin{circuitikz}[scale=0.9, transform shape]
				\draw
				% Fonte de 5V (positivo em cima, negativo no terra)
				(2.75,3.5) node[vcc](vcc){5V} (2.75,3.5)
				
				% Barramento positivo (topo)
				(2,3.5) -- (3.5,3.5)
				
				% Resistor 1 do paralelo (ramo esquerdo)
				(2,3.5) to[R, l=680$\Omega$] (2,1.75)
				
				
				% Resistor 2 do paralelo (ramo direito)
				
				
				(3.5,3.5) to[R, l=680$\Omega$] (3.5,1.75)
				
				% Barramento negativo do paralelo (fundo)
				(2,1.75) -- (3.5,1.75) -- (2.75,1.75)
				
				(2.75,1.75) to[ammeter] (2.75,0)
				% Resistor de 100Ω em série descendo para o terra
				-- (2.75,-0.3) to[R, l=100$\Omega$] (2.75,-1.5)
				
				% Terminal final conectado ao terra
				-- (2.75,-1.5) node[ground]{};
			\end{circuitikz}
		\end{center}
		
		O valor medido pelo multímetro de bancada foi de \boxed{12\ mA}
		
		\item \textbf{Tensão no resistor de 100 $\Omega$:}
		
		Para medir a tensão no resistor de 100 $\Omega$, foi colocado um multímetro em modo de medidor de tensão paralelamente ao resistor de 100 $\Omega$, de forma que o circuito resultante fique da seguinte forma:
		
		\begin{center}
			\begin{circuitikz}[scale=0.9, transform shape]
				\draw
				% Fonte de 5V (positivo em cima, negativo no terra)
				(2.75,3.5) node[vcc](vcc){5V} 
				
				% Barramento positivo (topo)
				(2,3.5) -- (3.5,3.5)
				
				% Resistor 1 do paralelo (ramo esquerdo)
				(2,3.5) to[R, l=680$\Omega$] (2,1.75)
				
				% Resistor 2 do paralelo (ramo direito)
				(3.5,3.5) to[R, l=680$\Omega$] (3.5,1.75)
				
				% Barramento negativo do paralelo (fundo)
				(2,1.75) -- (3.5,1.75) -- (2.75,1.75)
				
				% Resistor de 100Ω em série descendo para o terra
				(2.75,1.75) -- (2.75,1.5) 
				to[R, l=100$\Omega$, *-] (2.75,0)
				
				% Conectar o terra
				(2.75,0) node[ground]{}
				
				% ===============================================
				% VOLTÍMETRO EM PARALELO COM O RESISTOR DE 100Ω
				% ===============================================
				% O voltímetro vai do nó superior (2.75,1.5) ao nó inferior (2.75,0)
				% Mas para não sobrepor o resistor, colocamos ao lado esquerdo
				
				% Nó superior (antes do resistor)
				(2.75,1.5) -- (1.75,1.5)
				to[voltmeter, l=V] (1.75,0)
				-- (2.75,0)  % Nó inferior (terra)
				;
			\end{circuitikz}
		\end{center}
	\end{enumerate}
	
	O valor medido pelo multímetro de bancada foi de \boxed{1,11\ V}
	
\end{enumerate}


\section{Experimento 2}

\subsection{Montagem Experimental}

Foram utilizados 3 resistores de 1 k$\Omega$ com tolerância de $\pm$5\%. Os resistores foram posicionados na protoboard, o circuito não foi conectado à fontes de tensão, o circuito resultante foi o seguinte:

\begin{center}
\begin{circuitikz}[scale=0.9, transform shape]
	\draw
	
	(3,0) -- (3,-1) -- (0,-1) -- (0,0)
	
	to[R, l={1 k$\Omega$}] (1.5,0)
	to[R, l={1 k$\Omega$}, *-] (3,0)
	to[R, l={1 k$\Omega$}, *-] (4.5,0)
	
	(4.5,0) -- (4.5,1) -- (1.5,1) -- (1.5,0)
	;
\end{circuitikz}
\end{center}

\subsection{Resultados e Discussão}
\begin{enumerate}[label=\alph*]
	
	\item \textbf{Encontre a equação que descreve a associação de resistores presente na imagem.}
	
	Os 3 resistores presentes no circuito estão em paralelo entre sí, portanto podemos chegar na seguinte equação:
	
	$$\frac{1}{R_{eq}} = \frac{1}{R}+\frac{1}{R}+\frac{1}{R} \Rightarrow \boxed{R_{eq} = \frac{R}{3}}$$
	
	\item \textbf{Monte o circuito e meça a resistência da associação.}
	
	Após montar o circuito, foi medida a resistência, colocando o multímetro de bancada no modo de ohmímetro, a fonte de tensão contínua não foi conectada, a ponta positiva do multímetro foi colocada no ponto A e a negativa no ponto B, da seguinte forma:
	
	\begin{center}
		\begin{circuitikz}[scale=0.9, transform shape]
			\draw
			
			(3,0) -- (3,-1) -- (1.5, -1) node[circ, name=A, label=above:A]{} -- (0,-1) -- (0,0) 
			
			to[R, l={1 k$\Omega$}] (1.5,0)
			to[R, l={1 k$\Omega$}, *-] (3,0)
			to[R, l={1 k$\Omega$}, *-] (4.5,0)
			
			(4.5,0) -- (4.5,1) -- (3, 1) node[circ, name=A, label=above:B]{} -- (1.5,1) -- (1.5,0)
			;
		\end{circuitikz}
	\end{center}
	
	O valor medido pelo multímetro foi de \boxed{328\ \Omega}
	
\end{enumerate}


\section{Experimento 3}

\subsection{Montagem Experimental}

O circuito a seguir foi montado. A tensão de operação do LED representado pelo diodo D1 é de 3 V e a potência do mesmo é de 30 mW.
\begin{center}
	
	
	\begin{circuitikz}[scale=0.9, transform shape]
		\draw
		
		(-3,0) node[vcc](vcc){+5V} -- (-2,0)
		to[R, l=R1] (-0.5,0)
		-- (0.5,0)
		to[D, l=D1] (1.5,0)
		-- (2,0)
		node[ground]{}
		
		;
	\end{circuitikz}
\end{center}
\subsection{Resultados e Discussão}

\begin{enumerate}[label=\alph*]
	\item \textbf{Calcule o valor ideal para o resistor R1 e monte o circuito. Qual o valor real escolhido? Maior ou menor do que o ideal?}
	
	\begin{enumerate}
		\item \textbf {Valor ideal:}
		
		Dado que a tensão de operação do LED é de 3 V, e o valor fornecido pela fonte de tensão é de 5 V. Calcula-se então que a queda de tensão necessária no resistor R1 é de 2 V.
		Utilizando a lei de Watt, é possível encontrar a corrente que passa no resistor R1 e no LED:
		$$I_{LED} = I_{R1} = \frac{30\ mW}{3\ V} \Rightarrow I_{R1} = 10\ mA$$
		Utilizando a primeira lei de Ohm com o valor de corrente obtido, encontra-se o valor ideal de resistência:
		$$R1 = \frac{2\ V}{10\ mA} \Rightarrow \boxed{R1 = 200\ \Omega}$$
	
	
	\item \textbf{Qual valor real foi escolhido? Maior ou menor do que o ideal?}
	
	O valor real escolhido foi de $\boxed{330\ \Omega}$. Maior do que o valor ideal para garantir que não haja sobrecarga no diodo.
	\end{enumerate}
	
	\item \textbf{Construa o circuito com um potenciômetro. Qual o efeito da variação de resistência?}
	
	O seguinte circuito foi construído com um potenciômetro de 10 k$\Omega$.

	\begin{center}
		
		
		\begin{circuitikz}[scale=0.9, transform shape]
			\draw
			
			(-3,0) node[vcc](vcc){+5V} -- (-2.5,0)
			to[pR, l=$R_p$, name=POT] (-1,0)
			to[R, l=R1] (0.5,0)
			-- (0.5,0)
			to[D, l=D1] (1.5,0)
			-- (2,0)
			node[ground]{}
			
			;
		\end{circuitikz}
	\end{center}
	
	Conforme o seletor do potenciômetro é variado, há uma váriação na luminosidade do LED, com sua potência variando de forma linear com a resistência.
	
\end{enumerate}

\section{Dificuldades Encontradas}

A resistência real utilizada no experimento 3 foi muito alta, para aprimorar os testes, seria interessante utilizar uma resistência menor, mais próxima de 200 $\Omega$. Uma ótima opção seria um resistor de 220 $\Omega$, para garantir que o erro não seja um problema e que não haja sobrecarga no LED.
%\section{Anexos(Caso se aplique)}

%Caso o relatório deva ser acompanhado de outros arquivos (que serão também enviados via Moodle) indique aqui uma lista destes. Por ex:

%\begin{itemize}
%    \item Arquivo .c contendo determinado código.
%    \item Documentação de algum componente
%    \item etc
%\end{itemize}

\end{document}